\documentclass[11pt]{article}
\pdfinfoomitdate=1
\pdftrailerid{}
\pdfsuppressptexinfo=-1

\usepackage[a4paper,margin=1in]{geometry}
\usepackage{amsmath,amssymb,amsthm,mathtools}
\usepackage{booktabs}
\usepackage{enumitem}
\usepackage{microtype}
\usepackage[hidelinks]{hyperref}

\newtheorem{theorem}{Theorem}[section]
\newtheorem{proposition}[theorem]{Proposition}
\newtheorem{lemma}[theorem]{Lemma}
\newtheorem{corollary}[theorem]{Corollary}
\theoremstyle{definition}
\newtheorem{definition}[theorem]{Definition}
\newtheorem{remark}[theorem]{Remark}

\newtheorem*{theoremA}{Theorem A}
\newtheorem*{theoremB}{Theorem B}
\newtheorem*{theoremC}{Theorem C}
\newtheorem*{theoremD}{Theorem D}

\newcommand{\one}{\mathbf 1}
\newcommand{\supp}{\operatorname{supp}}
\newcommand{\degm}{\operatorname{deg}}
\newcommand{\R}{\mathbb R}
\newcommand{\Z}{\mathbb Z}
\newcommand{\E}{\mathbb E}

\title{Nested Regular Covers and Stopping Trees:\\
Finite Stabilization, Sharp Thresholds, and Integrality Gaps}
\author{}
\date{Private draft v1 --- 1 October 2026}

\begin{document}
\maketitle

\begin{abstract}
We study a nested covering problem for nonnegative regular matrices under dyadic capacity dilation.
A fixed $0$--$1$ mask $Z$ prescribes edges that must remain present, while a legal sequence
of regular matrices $(H_q)_{q\ge 0}$ satisfies
\[
Z\le H_0\le U,
\qquad
Z\le H_{q+1}\le 2H_q .
\]
The degree sequence is charged by the discounted objective
\[
C_\infty(H)=\sum_{q\ge0}2^{-q}\degm(H_q).
\]
This simple recurrence captures a temporal compatibility constraint that is invisible if each layer
is optimized independently.

We prove four results. First, if $r=|Z|$, every legal integer tail has a no-more-expensive representative
that becomes constant after
\[
M_r
=
\sum_{d=2}^{2r-1}
\left(1+\left\lceil\log_2(d-1)\right\rceil\right)
=
O(r\log r)
\]
layers, independently of the ambient dimension.
Second, for four required edges forming a partial matching, a local degree drop
$\kappa_Z(U)=3$ to $\kappa_Z(2U)=2$ requires at least seven states, and the bound is sharp.
Third, an explicit ten-state partial-matching instance has exact infinite-horizon optima
\[
C_\infty^{\Z}=6,
\qquad
C_\infty^{\R}=\frac{11}{2},
\]
so the integer/real gap is exactly $1/2$.
Finally, we give a closed-form stopping-tree realization for a six-state dyadic family with
\[
\mathbb E T_\ell
=
\frac{29}{8}-\frac38\,2^{-\ell},
\qquad
\mathbb E T_\infty=\frac{29}{8}.
\]

The general proofs are analytic.  The accompanying reproducibility package independently checks
the finite witnesses, the reachable-state certificate, the exact rational primal/dual certificate,
the constructive rounding mechanism, and the stopping-tree identities using exact arithmetic.
\end{abstract}

\section{Introduction}

Many exact randomized constructions are easy to optimize at a single time scale and harder to
optimize \emph{consistently across time}.  A local optimum at one resolution can force a support
pattern that obstructs a better optimum at the next resolution.  The present paper isolates a
finite combinatorial model of this phenomenon.

The state at a layer is a nonnegative regular matrix.  A lower mask $Z$ specifies edges that must
survive, while dyadic refinement doubles the available capacity.  The nesting condition
\[
H_{q+1}\le 2H_q
\]
forces every positive entry at a later layer to descend from positive mass already present at the
previous layer.  Consequently, support can shrink but cannot reappear.

Two questions motivate the paper.

\begin{enumerate}[label=(\roman*)]
\item Does the infinite-horizon optimization reduce to a finite computation whose size depends
only on the number $r=|Z|$ of required edges?
\item When do pointwise degree optima fail to fit into one common nested chain?
\end{enumerate}

Our first theorem answers the first question affirmatively with a state-count-independent
$O(r\log r)$ horizon.  The second and third theorems identify small sharp obstructions:
a four-edge partial matching already permits a local $3\to2$ capacity drop, but only from seven
states onward, while a five-edge ten-state construction separates the integer nested optimum from
its real relaxation by exactly $1/2$.

The final theorem reconnects the matrix problem to exact fair-bit stopping trees.  For an explicit
six-state family, cumulative stopped-branch matrices admit a closed form at every depth.  This
yields a complete finite tree for every member of the family and an infinite limiting tree with
finite expected cost.

\subsection{Classical background and related work}

The ingredients used here include several classical facts that are not novelty claims of this
paper.  Integer regular bipartite multigraphs decompose into perfect matchings, equivalently into
permutation matrices; this is the integral form behind the Birkhoff--von Neumann theorem
\cite{birkhoff1946,vonneumann1953}.  Lower/upper-capacity feasibility is naturally expressed
through integral flows and cut conditions of Hoffman type \cite{hoffman1974,bangjensengutin2009}.
Exact fair-bit sampling can be represented by binary discrete-distribution-generating trees in
the sense of Knuth and Yao \cite{knuthyao1976}; modern work continues to study exact and
entropy-efficient discrete sampling under random-bit models \cite{saad2020optimal}.

There is also a substantial literature on \emph{multistage} combinatorial optimization, where a
matching or matroid solution is selected at each time step and changing the selected solution
incurs acquisition or transition cost.  Gupta, Talwar, and Wieder introduced a general
multistage framework and showed strong hardness for perfect-matching maintenance
\cite{gupta2014changing}.  Bampis et al.\ studied multistage matching with explicit transition
costs \cite{bampis2018multistage}, and Chimani, Troost, and Wiedera developed approximation and
LP-integrality-gap results for intersection/union variants \cite{chimani2022approximating}.

The model here is different.  There is no exogenous sequence of edge costs and no penalty for
changing one matching into another.  Instead, the state is an integer regular \emph{multigraph}
matrix and temporal compatibility is imposed directly by the one-sided capacity recurrence
\[
H_{q+1}\le 2H_q .
\]
Thus later mass must descend from earlier support, while the objective discounts the regular
degree itself.  The sharp state-count threshold, the exact $6$ versus $11/2$ gap, and the
state-count-independent stabilization theorem below concern this nested-capacity model rather
than the standard multistage matching objectives.

Our contribution is the \emph{nested} interaction of regular covers under repeated dyadic
capacity dilation, together with the sharp finite obstructions and exact stopping certificates
developed below.  We deliberately separate analytic proofs from computational evidence.  The
supplied checker programs verify finite claims and certificates; they are not formal proofs of
the universal analytic theorems.

\subsection{Main results}

The four publication-level results are as follows.

\begin{theoremA}[Finite-layer stabilization]
Let $Z$ contain $r\ge2$ required edges.  Every legal integer tail has a no-more-expensive legal tail
that is constant after at most
\[
M_r
=
\sum_{d=2}^{2r-1}
\left(1+\left\lceil\log_2(d-1)\right\rceil\right)
\]
layers.  In particular $M_r=O(r\log r)$ and is independent of the matrix dimension.
\end{theoremA}

\begin{theoremB}[Sharp four-edge partial-matching threshold]
Let $Z$ consist of exactly four required edges in four distinct rows and four distinct columns.
If
\[
\kappa_Z(U)=3,
\qquad
\kappa_Z(2U)=2,
\]
then the matrix dimension satisfies $n\ge7$.  A seven-state example exists, so the bound is sharp.
\end{theoremB}

\begin{theoremC}[Exact integer/real integrality gap]
For the explicit ten-state partial-matching instance in Section~\ref{sec:gap},
\[
C_\infty^{\Z}=6,
\qquad
C_\infty^{\R}=\frac{11}{2}.
\]
Hence the exact integrality gap is $1/2$.
\end{theoremC}

\begin{theoremD}[Exact stopping-tree family]
For the six-state dyadic family in Section~\ref{sec:stopping},
there is an exact finite stopping tree for every $\ell\ge0$ with
\[
\mathbb E T_\ell
=
\frac{29}{8}-\frac38\,2^{-\ell}.
\]
The limiting kernel admits an exact infinite stopping tree with
\[
\mathbb E T_\infty=\frac{29}{8}.
\]
\end{theoremD}

\section{Regular covers and nested tails}
\label{sec:model}

\begin{definition}[Regular matrix]
A nonnegative $n\times n$ matrix $H$ is $d$-regular if
\[
H\one=d\one,
\qquad
H^\top\one=d\one.
\]
We write $\degm(H)=d$.
\end{definition}

Unless stated otherwise, regular matrices in the integer problem have entries in $\Z_{\ge0}$.

\begin{definition}[Required mask and covering degree]
Let $U\in\Z_{\ge0}^{n\times n}$ be regular and let
$Z\in\{0,1\}^{n\times n}$ satisfy $Z\le U$ entrywise.  Define
\[
\kappa_Z(U)
=
\min\left\{
d:
\exists\ d\text{-regular }H\in\Z_{\ge0}^{n\times n},
\quad
Z\le H\le U
\right\}.
\]
\end{definition}

\begin{definition}[Legal integer tail]
A legal tail for $(U,Z)$ is a sequence
$(H_q)_{q\ge0}$ of integer regular matrices satisfying
\[
Z\le H_0\le U,
\qquad
Z\le H_{q+1}\le2H_q
\quad(q\ge0).
\]
Writing $d_q=\degm(H_q)$, its discounted cost is
\[
C_\infty(H)
=
\sum_{q\ge0}2^{-q}d_q.
\]
The integer optimum is denoted $C_\infty^\Z(U,Z)$.
\end{definition}

The relaxed problem permits nonnegative real regular matrices and has optimum
$C_\infty^\R(U,Z)$.

We write $E_{ij}$ for the standard matrix unit and $J_n$ for the $n\times n$ all-ones matrix.

\begin{remark}[Support monotonicity]
If $H_{q,ij}=0$, then $(H_{q+1})_{ij}=0$.  Thus
\[
\supp H_{q+1}\subseteq\supp H_q.
\]
This elementary observation is the temporal constraint responsible for the phenomena below.
\end{remark}

\subsection{A cut formula for one layer}

The single-layer quantity $\kappa_Z(U)$ is a standard lower/upper-capacity flow problem
\cite{hoffman1974}.

\begin{proposition}[Cut characterization]
\label{prop:cut}
For row set $I$ and column set $J$, let
\[
s=|I|+|J|-n.
\]
For $s>0$, write $Z(I,J)$ for the total required mass in $I\times J$ and
$U(I^c,J^c)$ for the capacity in the complementary block.  Then
\[
\kappa_Z(U)
=
\max_{I,J:\ s>0}
\left\lceil
\frac{Z(I,J)-U(I^c,J^c)}{s}
\right\rceil.
\]
\end{proposition}

\begin{proof}
Let $H$ be $d$-regular with $Z\le H\le U$.  Regularity gives
\[
H(I,J)-H(I^c,J^c)=d(|I|+|J|-n)=ds.
\]
Using $H(I,J)\ge Z(I,J)$ and $H(I^c,J^c)\le U(I^c,J^c)$ yields the stated
lower bound on $d$.

Conversely, write $H=Z+X$ and impose the residual row and column demands needed to make
$H$ $d$-regular.  This is a bipartite flow problem with integral capacities
$0\le X\le U-Z$.  The max-flow/min-cut condition reduces exactly to the displayed family of
inequalities.  Integrality of the network implies an integral feasible $X$ whenever all cuts
pass.
\end{proof}

\section{Finite-layer stabilization}
\label{sec:finite}

This section proves Theorem~A.

\subsection{Degree trimming}

\begin{lemma}[Degree trimming]
\label{lem:trim}
Let $r=|Z|$.  Every legal tail has a no-more-expensive representative in which
\[
\degm(H_q)\le 2r-1
\qquad(q\ge0).
\]
\end{lemma}

\begin{proof}
Suppose a legal state $W\ge Z$ has degree at least $2r$.
Because $W$ is an integer regular bipartite multigraph, it decomposes into permutation matrices.
For each of the $r$ required edges, choose one permutation from the decomposition containing that
edge.  After removing duplicate choices, the sum $R$ satisfies
\[
Z\le R\le W,
\qquad
\degm(R)\le r.
\]
Replacing $W$ and its entire future by the constant tail $R,R,\ldots$ is legal.
If $W=H_q$ with $q>0$, the original chain has $W\le2H_{q-1}$, hence
$R\le W\le2H_{q-1}$.  If $q=0$, then $R\le W\le U$.
Every subsequent transition is legal because $R\le2R$.

Measured from the layer containing $W$, the replacement costs at most
\[
\degm(R)\sum_{j\ge0}2^{-j}
=
2\degm(R)
\le2r,
\]
whereas the original current layer alone costs at least $2r$.
Thus degree at least $2r$ is never necessary in an optimal representative.
\end{proof}

\subsection{A logarithmic reachability lemma}

\begin{lemma}[Reachability by balanced rounding]
\label{lem:reach}
Let $A$ be integer $d$-regular and $K$ integer $k$-regular with
\[
d>k,
\qquad
Z\le A,K,
\qquad
\supp K\subseteq\supp A.
\]
Then $K$ can be reached from $A$ by a legal chain with at most
\[
L(d)
=
1+\left\lceil\log_2(d-1)\right\rceil
\]
transitions, while every intermediate matrix before $K$ has degree $d$.
\end{lemma}

\begin{proof}
Choose a permutation matrix $P\le A$ and define
\[
R_0=K+(d-k)P.
\]
Then $R_0$ is $d$-regular, contains $Z$, and has support contained in $\supp A$.

Assume $R_s$ has these properties and put
\[
S_s=A+R_s.
\]
Every row and column sum of $S_s$ equals $2d$.
Begin with the entrywise floor $\lfloor S_s/2\rfloor$.
The set of positions at which $S_s$ is odd forms a bipartite graph in which every row-vertex and
column-vertex has even degree.  Decompose this graph into even cycles and alternately round up and
down around each cycle.  The resulting matrix $R_{s+1}$ is integer $d$-regular and satisfies
\[
R_{s+1,ij}
\in
\left\{
\left\lfloor\frac{A_{ij}+R_{s,ij}}2\right\rfloor,
\left\lceil\frac{A_{ij}+R_{s,ij}}2\right\rceil
\right\}.
\]
Because $A,R_s\ge Z$, every required edge survives.  Because both matrices are supported inside
$\supp A$, so is $R_{s+1}$.  Moreover
\[
R_s\le2R_{s+1}.
\]

Let
\[
e_s=\max_{ij}(R_{s,ij}-A_{ij})_+.
\]
Then
\[
e_{s+1}\le\left\lceil\frac{e_s}{2}\right\rceil.
\]
Every positive entry of $A$ is at least one and every entry of a $d$-regular nonnegative matrix is
at most $d$, so $e_0\le d-1$.  Hence after
\[
t=\left\lceil\log_2(d-1)\right\rceil
\]
rounds we have $e_t\le1$, equivalently $R_t\le2A$.

If $d=2$, then $k\le1$ and support containment implies $K\le A$, so the direct transition
$A,K$ is legal and the claimed bound is one.

Assume now $d\ge3$, so $t\ge1$.  The chain
\[
A,\ R_t,\ R_{t-1},\ldots,R_1,\ K
\]
is legal.  The first step uses $R_t\le2A$; the reversed rounding steps use
$R_s\le2R_{s+1}$; and the last step follows from $K\le R_0\le2R_1$.
The number of transitions is at most $1+t$.
\end{proof}

\subsection{Finite stabilization}

\begin{theorem}[Finite-layer stabilization]
\label{thm:finite}
Let $r=|Z|\ge2$ and set $B=2r-1$.  Every legal integer tail admits a no-more-expensive legal
tail that is constant after at most
\[
M_r
=
\sum_{d=2}^{B}
\left(1+\left\lceil\log_2(d-1)\right\rceil\right)
\]
layers.
If
\[
h_r=\left\lceil\log_2(B-1)\right\rceil,
\]
then
\[
M_r=(h_r+1)(B-1)-2^{h_r}+1,
\]
and therefore $M_r=O(r\log r)$.  In particular, the infinite-horizon infimum is attained.
\end{theorem}

\begin{proof}
Start from an arbitrary legal tail.  By Lemma~\ref{lem:trim}, replace it by a no-more-expensive
tail whose degrees are all at most $B$.  Along any legal tail, supports only shrink.

Now work with the running minimum of the degree.  Suppose $A$ is a state of current
running-minimum degree $d$, and let $K$ be the first later state whose degree is strictly smaller,
say $k<d$.  Every state between $A$ and $K$ has degree at least $d$, and support monotonicity gives
\[
\supp K\subseteq\supp A.
\]
Lemma~\ref{lem:reach} therefore replaces the whole segment by at most $L(d)$ degree-$d$ links
followed by $K$.

This replacement cannot increase discounted cost.  If $V_K$ denotes the original continuation
cost starting at $K$, then the constant tail $K,K,\ldots$ is feasible, so
\[
V_K\le2k<2d.
\]
Waiting $s$ degree-$d$ layers before entering that same continuation has value
\[
W_s
=
d\sum_{j=0}^{s-1}2^{-j}+2^{-s}V_K
=
2d+2^{-s}(V_K-2d),
\]
which is increasing in $s$.  Thus shortening the waiting segment to at most $L(d)$ layers cannot
increase cost.

If the degree never falls below the current running minimum again, freeze the current matrix
forever; this is legal and cannot increase cost because every later degree is at least the current
one.  Otherwise repeat the preceding acceleration at the next strict drop.  Each strict drop
lowers an integer degree, so this process terminates.

There are no relevant degrees above $B$.  Summing the worst-case delay over all possible strict
drops yields
\[
M_r=\sum_{d=2}^{B}L(d).
\]
The closed form follows by grouping the values of
$\lceil\log_2(d-1)\rceil$ into dyadic intervals.

Finally, after degree trimming there are only finitely many integer regular matrices of degree at
most $B$ on the fixed $n\times n$ support.  The preceding transformation reduces every legal tail
to one of finitely many tails of length at most $M_r$ followed by a constant state.  Hence the
infimum is a minimum.
\end{proof}

\section{The sharp four-edge partial-matching threshold}
\label{sec:four}

We now identify a sharp state-count obstruction for a four-edge partial matching.

\begin{theorem}[Four-edge lower bound]
\label{thm:fourlower}
Let $U$ be a nonnegative integer regular $n\times n$ capacity matrix and let $Z\le U$ consist of
exactly four required edges, with all four rows distinct and all four columns distinct.
If
\[
\kappa_Z(U)=3,
\qquad
\kappa_Z(2U)=2,
\]
then $n\ge7$.
\end{theorem}

\begin{proof}
Since $\kappa_Z(U)=3$, Proposition~\ref{prop:cut} yields a cut with $s>0$ such that
\[
Z(I,J)-U(I^c,J^c)\ge2s+1.
\]
Write
\[
z=Z(I,J),
\qquad
u=U(I^c,J^c).
\]
Because $\kappa_Z(2U)=2$, the same cut satisfies
\[
z-2u\le2s.
\]
Since $Z$ has four edges, $0\le z\le4$.  Thus
\[
z-u\ge2s+1,
\qquad
z-2u\le2s,
\qquad
0\le z\le4.
\]
The only integer solution is
\[
s=1,\qquad z=4,\qquad u=1.
\]
All four required edges therefore lie in $I\times J$.
Because they form a partial matching, $I$ contains four distinct required rows and $J$ contains
four distinct required columns.  Hence
\[
|I|\ge4,\qquad |J|\ge4.
\]
As $s=|I|+|J|-n=1$,
\[
n=|I|+|J|-1\ge7.
\]
\end{proof}

\begin{theorem}[Sharp seven-state witness]
\label{thm:fourseven}
The bound in Theorem~\ref{thm:fourlower} is sharp.
\end{theorem}

\begin{proof}
Let
\[
Z=E_{11}+E_{22}+E_{33}+E_{44}
\]
and
\[
U=
\begin{pmatrix}
1&0&0&0&1&1&0\\
0&1&0&0&1&0&1\\
0&0&1&0&0&1&1\\
0&0&0&1&0&1&1\\
1&1&0&0&1&0&0\\
1&0&1&1&0&0&0\\
0&1&1&1&0&0&0
\end{pmatrix}.
\]
The matrix $U$ is $3$-regular.  For $I=J=\{1,2,3,4\}$,
\[
s=1,\qquad Z(I,J)=4,\qquad U(I^c,J^c)=1.
\]
Proposition~\ref{prop:cut} therefore gives $\kappa_Z(U)\ge3$, and $U$ itself is a feasible
degree-three cover.

For $2U$, define
\[
H_2=
\begin{pmatrix}
1&0&0&0&0&1&0\\
0&1&0&0&0&0&1\\
0&0&1&0&0&1&0\\
0&0&0&1&0&0&1\\
0&0&0&0&2&0&0\\
1&0&1&0&0&0&0\\
0&1&0&1&0&0&0
\end{pmatrix}.
\]
Then $H_2$ is $2$-regular and $Z\le H_2\le2U$.
The same cut gives a lower bound of two, so $\kappa_Z(2U)=2$.
\end{proof}

\begin{remark}
The partial-matching assumption is essential to the state-count lower bound.
The theorem does not claim seven-state minimality for arbitrary four-edge masks with repeated
rows or columns.
\end{remark}

\section{An exact integer/real integrality gap}
\label{sec:gap}

We now pass from a local capacity drop to a genuinely infinite nested optimization.
The theorem is stated directly from the following ten-state partial-matching instance; no
auxiliary construction is needed for its proof.

Let $U'\in\Z_{\ge0}^{10\times10}$ be
{\small
\[
U'=
\begin{pmatrix}
1&0&1&0&0&2&0&0&0&0\\
0&1&0&0&2&0&1&0&0&0\\
0&0&1&0&1&1&0&0&1&0\\
0&2&0&2&0&0&0&0&0&0\\
0&0&2&1&1&0&0&0&0&0\\
2&0&0&1&0&1&0&0&0&0\\
0&0&0&0&0&0&3&1&0&0\\
1&0&0&0&0&0&0&3&0&0\\
0&0&0&0&0&0&0&0&3&1\\
0&1&0&0&0&0&0&0&0&3
\end{pmatrix}
\]
}
and let
\[
Z'=E_{11}+E_{22}+E_{33}+E_{78}+E_{9,10}.
\]

\subsection{A persistent degree-two lower bound}

\begin{lemma}[All tail layers have degree at least two]
\label{lem:tail2}
For every legal integer or real tail for $(U',Z')$,
\[
d_q\ge2
\qquad(q\ge0).
\]
\end{lemma}

\begin{proof}
Support monotonicity and $H_0\le U'$ imply
\[
\supp H_q\subseteq\supp U'
\qquad(q\ge0).
\]
Row $8$ of $U'$ is supported only on columns $1$ and $8$.
The required mask has the edge $(1,1)$ in column $1$ and the edge $(7,8)$ in column $8$, both
outside row $8$.
For a $d_q$-regular feasible layer, the total mass in columns $1$ and $8$ is $2d_q$.
Row $8$ contributes exactly $d_q$ to those columns, while the two required edges outside row $8$
contribute at least two more.  Hence
\[
2d_q\ge d_q+2,
\]
and therefore $d_q\ge2$.
\end{proof}

\subsection{Integer optimum}

\begin{lemma}[Reachable degree-three graph]
\label{lem:reachgraph}
There is no integer entrance state of degree at most two, and there is a unique integer
degree-three entrance state
\[
Z'\le H_0\le U'.
\]
Under degree-three transitions $H'\le2H$, exactly four degree-three states are reachable, with
nine directed transitions in total, and none of the four admits a degree-two successor containing
$Z'$.
\end{lemma}

\begin{proof}
This is a finite exact enumeration of regular matrices under the displayed capacities.
The full state list and transitions are certified by the accompanying checker.
No floating-point optimization is used.
\end{proof}

\begin{proposition}[Integer optimum]
\label{prop:intopt}
For $(U',Z')$,
\[
C_\infty^\Z(U',Z')=6.
\]
\end{proposition}

\begin{proof}
A constant degree-three chain is feasible, hence
\[
C_\infty^\Z\le3\sum_{q\ge0}2^{-q}=6.
\]

For the reverse inequality, Lemma~\ref{lem:reachgraph} gives $d_0\ge3$.
If a legal integer chain never reaches degree two, then every layer has degree at least three and
its cost is at least six.

Otherwise let $t$ be the first degree-two layer.  Before time $t$, all degrees are at least three.
By Lemma~\ref{lem:reachgraph}, a degree-three state in the reachable sector has no degree-two
successor, so at least one layer $j<t$ has degree at least four.
Relative to the constant degree-three chain, that layer contributes at least $2^{-j}$ extra.
By Lemma~\ref{lem:tail2}, no later layer can have degree below two, so the maximum possible total
saving after time $j$ is
\[
\sum_{q\ge j+1}2^{-q}=2^{-j}.
\]
Thus the increase cannot yield a net improvement, and every legal integer chain has cost at least
six.
\end{proof}

\subsection{Real optimum}

The following two matrices form an explicit relaxed entrance and tail:
{\scriptsize
\[
H_0^\R=
\begin{pmatrix}
1&0&1/2&0&0&2&0&0&0&0\\
0&1&0&0&3/2&0&1&0&0&0\\
0&0&1&0&1&1/2&0&0&1&0\\
0&3/2&0&2&0&0&0&0&0&0\\
0&0&2&1/2&1&0&0&0&0&0\\
3/2&0&0&1&0&1&0&0&0&0\\
0&0&0&0&0&0&5/2&1&0&0\\
1&0&0&0&0&0&0&5/2&0&0\\
0&0&0&0&0&0&0&0&5/2&1\\
0&1&0&0&0&0&0&0&0&5/2
\end{pmatrix}.
\]
\[
H_1^\R=
\begin{pmatrix}
1&0&1&0&0&0&0&0&0&0\\
0&1&0&0&0&0&1&0&0&0\\
0&0&1&0&0&0&0&0&1&0\\
0&0&0&2&0&0&0&0&0&0\\
0&0&0&0&2&0&0&0&0&0\\
0&0&0&0&0&2&0&0&0&0\\
0&0&0&0&0&0&1&1&0&0\\
1&0&0&0&0&0&0&1&0&0\\
0&0&0&0&0&0&0&0&1&1\\
0&1&0&0&0&0&0&0&0&1
\end{pmatrix}.
\]
}
Their degrees are $7/2$ and $2$ respectively, and
\[
Z'\le H_0^\R\le U',
\qquad
Z'\le H_1^\R\le2H_0^\R.
\]
Taking $H_q^\R=H_1^\R$ for all $q\ge1$ gives
\[
C_\infty^\R\le\frac72+2=\frac{11}{2}.
\]

\begin{lemma}[Exact dual inequality]
\label{lem:dual}
Every feasible real pair $(H_0,H_1)$ for the first two layers satisfies
\[
2d_0+d_1\ge9.
\]
\end{lemma}

\begin{proof}
The inequality has an exact rational linear-programming dual certificate.
It is obtained as an integer linear combination of row-regularity equalities,
column-regularity equalities, one nesting inequality, active lower bounds, and active upper
capacity bounds.  The complete coefficient list is reproduced in Appendix~\ref{app:dual} and is
checked symbolically by the reproducibility package.
\end{proof}

\begin{proposition}[Real optimum]
\label{prop:realopt}
For $(U',Z')$,
\[
C_\infty^\R(U',Z')=\frac{11}{2}.
\]
\end{proposition}

\begin{proof}
The explicit chain above gives the upper bound.
For the lower bound, Lemmas~\ref{lem:dual} and~\ref{lem:tail2} imply
\[
\begin{aligned}
C_\infty^\R
&=
d_0+\frac12d_1+\sum_{q\ge2}2^{-q}d_q\\
&\ge
\frac12(2d_0+d_1)
+
2\sum_{q\ge2}2^{-q}\\
&\ge
\frac92+1
=
\frac{11}{2}.
\end{aligned}
\]
\end{proof}

Combining Propositions~\ref{prop:intopt} and~\ref{prop:realopt} proves Theorem~C.

\section{Stopping-tree realization}
\label{sec:stopping}

We now connect the same nested-cover mechanism to exact stopping trees through a separate
six-state five-edge construction.  A small zero-sum correction of its negative support makes the
matrix family bistochastic at every dyadic scale, allowing nested cover data to be read as
cumulative stopping data.

\subsection{Dyadic kernels and stopping trees}

A matrix $A$ is \emph{bistochastic} if it is entrywise nonnegative and
\[
A\one=\one,
\qquad
A^\top\one=\one.
\]
It is \emph{dyadic} if every entry is a dyadic rational.  Its dyadic height is the least
$d\ge0$ for which $2^dA$ is an integer matrix.  Floors of matrices are taken entrywise.

A binary fair-bit stopping tree has probability $2^{-t}$ on each node at depth $t$.
A terminal node is labeled by a permutation matrix.  The tree \emph{realizes} $A$ if it
terminates almost surely and the probability-weighted sum of its terminal permutation matrices
is exactly $A$.

\subsection{Cumulative stopping matrices}

Consider such a tree.
For depth $t$, let $N_t$ be the sum of the permutation matrices attached to all leaves that have
stopped by time $t$, counted at depth-$t$ multiplicity.
If $N_t$ is $s_t$-regular, then the probability of still running after $t$ bits is
\[
1-\frac{s_t}{2^t}.
\]

\begin{proposition}[Integer nesting criterion]
\label{prop:nesting}
Let $A$ be a dyadic bistochastic matrix of height $d$.
A cumulative stopping profile is realizable by a permutation-valued binary tree of depth $d$ if
there are integer regular matrices $N_0,\ldots,N_d$ such that
\[
0\le N_t\le\lfloor2^tA\rfloor,
\qquad
N_{t+1}\ge2N_t,
\qquad
N_d=2^dA.
\]
\end{proposition}

\begin{proof}
Necessity is immediate from the tree.
For sufficiency, let $s_t=\degm(N_t)$ and define the new stopping matrix
\[
X_t=N_t-2N_{t-1}
\]
(with $N_{-1}=0$ and $s_{-1}=0$).
Then $X_t$ is nonnegative integer regular of degree
\[
x_t=s_t-2s_{t-1}.
\]
Decompose $X_t$ into $x_t$ permutation matrices.
Before stopping new nodes at depth $t$, the number of unresolved binary-tree nodes is
\[
2^t-2s_{t-1}.
\]
The capacity inequality implies $s_t\le2^t$, hence
\[
x_t=s_t-2s_{t-1}\le2^t-2s_{t-1}.
\]
We may therefore assign the $x_t$ permutation matrices to distinct unresolved nodes and leave the
remaining nodes alive.  Induction over $t$ realizes every $N_t$.
Finally, $N_d=2^dA$ has degree $2^d$, so no live nodes remain at depth $d$ and the resulting tree
realizes $A$ exactly.
\end{proof}

\begin{corollary}[Infinite nesting criterion]
\label{cor:infinite-nesting}
Let $A$ be bistochastic.  Suppose $(N_t)_{t\ge0}$ is a sequence of nonnegative integer regular
matrices with degrees $s_t$ such that
\[
N_{t+1}\ge2N_t,
\qquad
N_t\le\lfloor2^tA\rfloor
\quad(t\ge0),
\]
and
\[
\frac{N_t}{2^t}\longrightarrow A
\]
entrywise.  Then there is a permutation-valued infinite binary stopping tree realizing $A$, and
it terminates almost surely.
\end{corollary}

\begin{proof}
Set $X_t=N_t-2N_{t-1}$ as in Proposition~\ref{prop:nesting}.  Each $X_t$ is a nonnegative integer
regular matrix of degree $x_t=s_t-2s_{t-1}$ and therefore decomposes into permutation matrices.
Since $N_t\le\lfloor2^tA\rfloor$, we have $s_t\le2^t$, so
\[
x_t\le2^t-2s_{t-1},
\]
the number of unresolved nodes available at depth $t$.  We may therefore assign the permutation
summands of $X_t$ to new stopping nodes inductively at every finite depth.

At depth $t$, the cumulative realized matrix is $N_t/2^t$, which converges entrywise to $A$.
Taking row sums and using bistochasticity gives
\[
\frac{s_t}{2^t}\longrightarrow1.
\]
Hence the unresolved probability
\[
1-\frac{s_t}{2^t}
\]
tends to zero, so the tree terminates almost surely and its terminal distribution realizes $A$.
\end{proof}

\subsection{An explicit six-state family}

Let
\[
U=
\begin{pmatrix}
1&0&1&0&0&2\\
1&1&0&0&2&0\\
0&1&1&0&1&1\\
0&2&0&2&0&0\\
0&0&2&1&1&0\\
2&0&0&1&0&1
\end{pmatrix}.
\]
Let $P$ be the permutation matrix with row-columns $(3,5,6,2,4,1)$ and define
\[
H=U-P.
\]
The original five-edge core mask is
\[
Z_{\rm core}
=
E_{11}+E_{21}+E_{22}+E_{32}+E_{33}.
\]
For the stochastic embedding, use the corrected negative support
\[
\widehat Z
=
Z_{\rm core}+E_{44}+E_{55},
\]
and
\[
\widehat G
=
E_{12}+E_{23}+E_{24}+E_{31}+E_{35}+E_{41}+E_{52}
-
\widehat Z.
\]
Then
\[
\widehat G\one=0,
\qquad
\widehat G^\top\one=0,
\]
and $H$ is $3$-regular with $H\ge\widehat Z$.

Define
\[
A_\infty=\frac{2J_6+U}{16},
\qquad
B=16A_\infty=2J_6+U,
\]
and for $\ell\ge0$,
\[
A_{4+\ell}
=
A_\infty+2^{-(4+\ell)}\widehat G.
\]
Set $s=4+\ell$.
Because $U$ is $4$-regular and $\widehat G$ has zero row and column sums, every $A_s$ is
bistochastic.  Moreover $U+\widehat G\ge0$, so $A_4\ge0$; for $\ell>0$ the perturbation has
smaller magnitude entrywise, hence $A_s\ge0$ as well.  Thus every $A_s$ is a valid dyadic
bistochastic kernel.

\begin{theorem}[Finite stopping-tree family]
\label{thm:finite-tree}
For each $\ell\ge0$, define cumulative matrices by
\[
N_0=N_1=N_2=0,
\qquad
N_3=J_6,
\]
\[
N_t=2^tA_\infty-H=2^{t-4}B-H
\qquad(4\le t<s),
\]
and
\[
N_s=2^sA_s=2^{s-4}B+\widehat G.
\]
These matrices satisfy Proposition~\ref{prop:nesting}; hence they define an exact stopping tree for
$A_s$.  Its expected fair-bit cost is
\[
\mathbb E T_\ell
=
\frac{29}{8}-\frac38\,2^{-\ell}.
\]
\end{theorem}

\begin{proof}
Regularity is immediate from the row and column sums.
For $t\le2$ the capacity bound is trivial because $N_t=0$.
At $t=3$,
\[
8A_s=\frac{B}{2}+2^{3-s}\widehat G.
\]
Every entry of $B$ is at least two.  On the negative support of $\widehat G$, the corresponding
entry of $U$ is at least one, so the corresponding entry of $B$ is at least three.  Since
$0<2^{3-s}\le1/2$, every entry of $8A_s$ is at least one.  Hence
\[
N_3=J_6\le\lfloor8A_s\rfloor.
\]

For $4\le t<s$, entrywise flooring of $2^tA_s$ can reduce an entry below
$2^{t-4}B$ only on the negative support of $\widehat G$.
Since $H\ge\widehat Z$, we have
\[
N_t\le\lfloor2^tA_s\rfloor.
\]
At the terminal depth, equality $N_s=2^sA_s$ holds by definition.

The nesting increments are
\[
N_4-2N_3=P\ge0,
\]
at a nonterminal level four,
\[
N_{t+1}-2N_t=H\ge0
\]
between intermediate levels, and
\[
N_s-2N_{s-1}=\widehat G+2H\ge0
\]
at termination.  For $\ell=0$, the direct terminal increment is $U+\widehat G\ge0$.

The cumulative degrees are
\[
\degm N_t=
\begin{cases}
0,&t=0,1,2,\\
6,&t=3,\\
2^t-3,&4\le t<s,\\
2^s,&t=s.
\end{cases}
\]
Therefore
\[
\begin{aligned}
\mathbb E T_\ell
&=
\sum_{t=0}^{s-1}
\left(1-\frac{\degm N_t}{2^t}\right)\\
&=
3+\left(1-\frac68\right)+\sum_{t=4}^{s-1}\frac{3}{2^t}\\
&=
\frac{29}{8}-\frac38\,2^{-\ell}.
\end{aligned}
\]
\end{proof}

\begin{theorem}[Infinite limiting tree]
\label{thm:infinite-tree}
The limiting kernel $A_\infty$ admits an exact infinite stopping tree with
\[
N_0=N_1=N_2=0,\qquad N_3=J_6,
\]
and
\[
N_t=2^tA_\infty-H
\qquad(t\ge4).
\]
Its expected cost is
\[
\mathbb E T_\infty=\frac{29}{8}.
\]
\end{theorem}

\begin{proof}
For $t\ge4$,
\[
N_{t+1}-2N_t=H\ge0.
\]
Moreover
\[
\frac{N_t}{2^t}
=
A_\infty-\frac{H}{2^t}
\longrightarrow A_\infty
\]
entrywise, while the surviving mass at level $t$ is exactly $3/2^t$.
Corollary~\ref{cor:infinite-nesting} therefore gives an exact infinite stopping tree, and the
surviving mass tends to zero, so the tree terminates almost surely.
Summing the survival probabilities gives
\[
3+\left(1-\frac68\right)+\sum_{t\ge4}\frac{3}{2^t}
=
\frac{29}{8}.
\]
\end{proof}

Theorems~\ref{thm:finite-tree} and~\ref{thm:infinite-tree} prove Theorem~D.

\section{Computational certificates and reproducibility}
\label{sec:repro}

The reproducibility package separates analytic proof from finite computation.

\begin{center}
\small
\begin{tabular}{@{}l p{0.31\textwidth} p{0.47\textwidth}@{}}
\toprule
Result & Finite checker & What is checked\\
\midrule
Theorem A &
\path{reachability_rounding.py} &
balanced-rounding construction and $M_r$ formula\\
Theorem B &
\path{four_edge_minimality.py} &
cut arithmetic and sharp seven-state witness\\
Theorem C &
\path{integrality_gap.py} &
integer state graph, real primal, and exact dual certificate\\
Theorem D &
\path{stopping_tree.py} &
symbolic local identities, finite family, and infinite tail\\
\bottomrule
\end{tabular}
\end{center}

All theorem-facing pass/fail arithmetic in these scripts is integer or rational.
A single public entry point runs the four suites and verifies the SHA-256 locked publication
snapshot.

The computation does not replace the analytic arguments in
Theorem~\ref{thm:finite} or Theorem~\ref{thm:fourlower}.
For Theorem~C, the finite state graph and rational dual identity are part of the certificate used
by the proof; for Theorem~D, the checker verifies the explicit symbolic matrix identities.

\section{Discussion}

The results show that temporal compatibility introduces structure not visible in independent
single-layer optimization.

First, the infinite problem is nevertheless finitely determined: after degree trimming, every
strict degree drop can be accelerated through a logarithmic number of balanced-rounding layers.
The resulting $O(r\log r)$ bound depends only on the number of required edges.

Second, local capacity improvement has a sharp small obstruction.
For a four-edge partial matching, a $3\to2$ drop cannot occur below seven states.

Third, the nested integer problem is genuinely nonconvex even when its real relaxation is simple.
The ten-state construction gives the exact gap
\[
6-\frac{11}{2}=\frac12.
\]

Finally, the six-state family shows that the matrix certificates are operational: they can be
assembled into complete exact stopping trees, including an infinite limiting tree of finite mean.

Several directions remain outside the present publication unit.  The broader classification of
simultaneous optimality for other support geometries, sharper bounds on the stabilization horizon,
and structural conditions guaranteeing equality of integer and real nested optima are natural
next questions.  These problems are intentionally not folded into the present theorem set.

\appendix

\section{Exact dual certificate for Theorem C}
\label{app:dual}

We record the coefficient data used to verify
\[
2d_0+d_1\ge9.
\]
Variables are the supported entries of the first two real layers together with $d_0,d_1$.
The certificate combines row and column regularity equalities, one nesting inequality, lower
bounds, and upper-capacity bounds.

The nonzero equality multipliers are:
\[
\begin{array}{c|r}
\text{constraint} & \text{multiplier}\\
\hline
(0,\mathrm{row}\ 4),(0,\mathrm{row}\ 5),(0,\mathrm{row}\ 6),
(0,\mathrm{row}\ 8),(0,\mathrm{row}\ 10) & 2\\
(0,\mathrm{col}\ 1),(0,\mathrm{col}\ 2),(0,\mathrm{col}\ 3),
(0,\mathrm{col}\ 4),(0,\mathrm{col}\ 8),(0,\mathrm{col}\ 10) & -2\\
(1,\mathrm{row}\ 2),(1,\mathrm{row}\ 3),(1,\mathrm{row}\ 5),
(1,\mathrm{row}\ 7),(1,\mathrm{row}\ 9) & -1\\
(1,\mathrm{col}\ 3),(1,\mathrm{col}\ 5),
(1,\mathrm{col}\ 7),(1,\mathrm{col}\ 9) & 1
\end{array}
\]
together with multiplier $-1$ on the nesting inequality at entry $(1,3)$ in one-based notation.

The positive lower-bound multipliers are $2$ on
\[
H_{0,11},\ H_{0,22},\ H_{0,33},\ H_{0,78},\ H_{0,9,10}\ge1,
\]
and $1$ on
\[
H_{1,22}\ge1,\quad
H_{1,36}\ge0,\quad
H_{1,54}\ge0,\quad
H_{1,78}\ge1,\quad
H_{1,9,10}\ge1.
\]
The negative upper-bound multipliers are $-2$ on
\[
H_{0,55}\le1,
\qquad
H_{0,66}\le1.
\]
After collecting coefficients, every matrix variable cancels and the left-hand side becomes
$2d_0+d_1$, while the bound constants sum to $9$.
The accompanying exact checker evaluates the complete coefficient vector independently.

\section{Reproducibility statement}

The publication snapshot is generated from a private canonical repository by an explicit
allowlist manifest.  Every exported artifact is SHA-256 pinned.  The exporter fails if any
canonical artifact drifts from its frozen digest.

The public command
\[
\texttt{python verify.py}
\]
checks snapshot integrity and independently runs all finite certificate suites.
Its intended interpretation is:

\begin{quote}
The command verifies the finite certificates and computational claims shipped with this
publication snapshot.  Analytic proofs are contained in the paper and are not formally verified
by these scripts.
\end{quote}

\bibliographystyle{plain}
\bibliography{references}

\end{document}
